% ---------------------------------------------------------------------
%  PHỤ LỤC E. DANH MỤC SƠ ĐỒ VÀ MÃ NGUỒN PLANTUML
% ---------------------------------------------------------------------
\section{Danh mục sơ đồ và mã nguồn PlantUML}\label{pl:so-do}

Các sơ đồ trong báo cáo được viết bằng PlantUML~\cite{plantuml}, lưu dưới dạng văn bản trong \texttt{diagrams/} và kết xuất thành ảnh PNG trong \texttt{figures/diagrams/}. Lưu mã nguồn giúp quản lý phiên bản và so sánh thay đổi; kiểu dáng chung được khai báo trong \texttt{\_style.iuml}. Bảng~\ref{tab:danh-muc-so-do} liệt kê sơ đồ và tệp nguồn tương ứng.

{\small
\begin{longtable}{|C{1.3cm}|L{\dimexpr\textwidth-8\tabcolsep-5\arrayrulewidth-1.3cm-5.6cm-2.4cm\relax}|L{5.6cm}|L{2.4cm}|}
\caption{Danh mục sơ đồ và tệp nguồn PlantUML}\label{tab:danh-muc-so-do}\\
\hline
\hd{Hình} & \hd{Tên sơ đồ} & \hd{Tệp nguồn} & \hd{Loại sơ đồ}\\ \hline
\endfirsthead
\multicolumn{4}{l}{\footnotesize\itshape Bảng \thetable\ (tiếp theo)}\\ \hline
\hd{Hình} & \hd{Tên sơ đồ} & \hd{Tệp nguồn} & \hd{Loại sơ đồ}\\ \hline
\endhead
\ref{fig:quy-trinh-thu-cong} & Quy trình đi chợ và sử dụng thực phẩm theo cách thủ công & \path{c1-01-quy-trinh-thu-cong.puml} & Hoạt động\\ \hline
\ref{fig:phan-ra-chuc-nang} & Sơ đồ phân rã chức năng của ứng dụng & \path{c1-03-phan-ra-chuc-nang.puml} & Sơ đồ tư duy\\ \hline
\ref{fig:chu-trinh} & Chu trình nghiệp vụ khép kín đề xuất cho ứng dụng & \path{c1-02-chu-trinh-nghiep-vu.puml} & Khối quy trình\\ \hline
\ref{fig:ngu-canh} & Sơ đồ ngữ cảnh của hệ thống & \path{c1-04-ngu-canh.puml} & Ngữ cảnh\\ \hline
\ref{fig:ipo-tong-quat} & Mô hình Đầu vào – Xử lý – Đầu ra tổng quát của hệ thống & \path{c1-05-ipo-tong-quat.puml} & Khối IPO\\ \hline
\ref{fig:quy-trinh-mua-sam} & Quy trình mua sắm có phân công trong nhóm gia đình & \path{c1-06-quy-trinh-mua-sam-nhom.puml} & Hoạt động\\ \hline
\ref{fig:trang-thai-mat-hang} & Vòng đời của một mặt hàng trong danh sách mua sắm & \path{c1-07-trang-thai-mat-hang.puml} & Trạng thái\\ \hline
\ref{fig:trang-thai-thuc-pham} & Vòng đời của một lô thực phẩm trong tủ lạnh & \path{c1-08-trang-thai-thuc-pham.puml} & Trạng thái\\ \hline
\ref{fig:quy-trinh-ke-hoach} & Quy trình lập kế hoạch bữa ăn gắn với tồn kho & \path{c1-09-quy-trinh-ke-hoach-bua-an.puml} & Hoạt động\\ \hline
\ref{fig:mo-hinh-khai-niem} & Mô hình khái niệm nghiệp vụ & \path{c1-10-mo-hinh-khai-niem.puml} & Lớp khái niệm\\ \hline
\ref{fig:quan-he-tac-nhan} & Biểu đồ quan hệ giữa các tác nhân & \path{c2-01-quan-he-tac-nhan.puml} & Use case\\ \hline
\ref{fig:uc-tong-quan} & Biểu đồ use case tổng quan của hệ thống & \path{c2-02-uc-tong-quan.puml} & Use case\\ \hline
\ref{fig:uc-tai-khoan} & Biểu đồ use case phân rã gói Xác thực và quản lý tài khoản & \path{c2-03-uc-tai-khoan.puml} & Use case\\ \hline
\ref{fig:act-dang-ky} & Biểu đồ hoạt động của use case Đăng ký tài khoản và Xác thực mã OTP & \path{c2-04-act-dang-ky.puml} & Hoạt động\\ \hline
\ref{fig:uc-nhom} & Biểu đồ use case phân rã gói Quản lý nhóm gia đình & \path{c2-05-uc-nhom.puml} & Use case\\ \hline
\ref{fig:act-tham-gia-nhom} & Biểu đồ hoạt động của use case Tham gia nhóm gia đình & \path{c2-06-act-tham-gia-nhom.puml} & Hoạt động\\ \hline
\ref{fig:uc-mua-sam} & Biểu đồ use case phân rã gói Quản lý danh sách mua sắm & \path{c2-07-uc-mua-sam.puml} & Use case\\ \hline
\ref{fig:act-danh-dau-da-mua} & Biểu đồ hoạt động của use case Đánh dấu đã mua và Chuyển vào tủ lạnh & \path{c2-08-act-danh-dau-da-mua.puml} & Hoạt động\\ \hline
\ref{fig:act-sinh-ds} & Biểu đồ hoạt động của use case Sinh danh sách mua sắm từ kế hoạch bữa ăn & \path{c2-09-act-sinh-ds-tu-ke-hoach.puml} & Hoạt động\\ \hline
\ref{fig:uc-tu-lanh} & Biểu đồ use case phân rã gói Quản lý thực phẩm trong tủ lạnh & \path{c2-10-uc-tu-lanh.puml} & Use case\\ \hline
\ref{fig:act-nhac-nho} & Biểu đồ hoạt động của use case Gửi nhắc nhở thực phẩm sắp hết hạn & \path{c2-11-act-nhac-nho.puml} & Hoạt động\\ \hline
\ref{fig:seq-nhac-nho} & Biểu đồ tuần tự của luồng nhắc nhở hạn dùng & \path{c2-12-seq-nhac-nho.puml} & Tuần tự\\ \hline
\ref{fig:uc-cong-thuc} & Biểu đồ use case phân rã gói Quản lý công thức nấu ăn & \path{c2-13-uc-cong-thuc.puml} & Use case\\ \hline
\ref{fig:uc-ke-hoach} & Biểu đồ use case phân rã gói Quản lý kế hoạch bữa ăn & \path{c2-14-uc-ke-hoach.puml} & Use case\\ \hline
\ref{fig:act-goi-y} & Biểu đồ hoạt động của use case Gợi ý món ăn từ thực phẩm sẵn có & \path{c2-15-act-goi-y-mon.puml} & Hoạt động\\ \hline
\ref{fig:act-xac-nhan-nau} & Biểu đồ hoạt động của use case Xác nhận đã nấu món ăn & \path{c2-16-act-xac-nhan-nau.puml} & Hoạt động\\ \hline
\ref{fig:uc-tim-kiem} & Biểu đồ use case phân rã gói Tìm kiếm và phân loại thực phẩm & \path{c2-17-uc-tim-kiem.puml} & Use case\\ \hline
\ref{fig:uc-thong-ke} & Biểu đồ use case phân rã gói Thống kê báo cáo & \path{c2-18-uc-thong-ke.puml} & Use case\\ \hline
\ref{fig:uc-quan-tri} & Biểu đồ use case phân rã gói Quản trị hệ thống & \path{c2-19-uc-quan-tri.puml} & Use case\\ \hline
\end{longtable}
}

Sau khi chỉnh sửa sơ đồ, chạy \texttt{make diagrams} tại thư mục gốc để kết xuất lại các tệp \texttt{diagrams/*.puml} thành ảnh PNG. Các sơ đồ use case, trạng thái, lớp và ngữ cảnh dùng bộ dàn trang Smetana tích hợp trong PlantUML, không cần Graphviz. Hướng dẫn chi tiết có trong \texttt{README.md}.

Tệp kiểu dáng dùng chung được trình bày dưới đây. Mỗi sơ đồ nạp tệp này bằng chỉ thị \texttt{!include \_style.iuml}.

\codecaption{Tệp kiểu dáng dùng chung \_style.iuml}{lst:style}
\VerbatimInput{diagrams/_style.iuml}

Mã nguồn dưới đây là biểu đồ phân rã gói Xác thực và quản lý tài khoản (Hình~\ref{fig:uc-tai-khoan}), minh họa cách khai báo tác nhân, use case, quan hệ bao gồm và tác nhân phụ.

\codecaption{Tệp c2-03-uc-tai-khoan.puml}{lst:uc-tai-khoan}
\VerbatimInput{diagrams/c2-03-uc-tai-khoan.puml}

Biểu đồ hoạt động Gợi ý món ăn (Hình~\ref{fig:act-goi-y}) minh họa cách mô tả làn, vòng lặp và nhánh xử lý trong PlantUML.

\codecaption{Tệp c2-15-act-goi-y-mon.puml}{lst:act-goi-y}
\VerbatimInput{diagrams/c2-15-act-goi-y-mon.puml}
